\section{Polynomial identities behind the archimedean coefficients}
\label{ca-section}

The next step concerns the differential operators that produce
the archimedean Fourier estimates.  We prove their leading
polynomial identity directly, including the two different orders
of the matrix indices, and then prove the coefficient extraction
and its positive matrix estimate.  We also bound the explicit
rational correction contributed by the preceding operators.
These are actual algebraic steps in the source's coefficient
argument~\cite{alpo-uniform-source}.  The Fourier integral
identity for the radial Hermitian form, and positivity of that
form, are identified separately at the end.

\subsection{The differential operators and their leading coefficients}

Let $k,r$ be nonnegative integers and $n=2r+k+1$.
Use the ordered basis
\[
 e_0,\ldots,e_{r-1},w_0,\ldots,w_k,f_{r-1},\ldots,f_0
\]
of $\mathbb C^n$.  Let $H$ be the Hermitian matrix specified by
\[
 He_a=f_a,\quad Hf_a=e_a,\quad Hw_b=w_b\ (b<k),
 \quad Hw_k=-w_k.
\]
Thus $H^*=H$ and $H^2=1$.  Put
\[
 G=\{g\in\operatorname{GL}_n(\mathbb C):g^*Hg=H\}.
\]
Its real Lie algebra consists of the matrices $X$ satisfying
$X^*H+HX=0$.

For a matrix $Z$ set
\[
 Z_1=\tfrac12(Z-HZ^*H),\qquad
 Z_2=\tfrac1{2i}(Z+HZ^*H).
\]
Both matrices belong to this real Lie algebra, and
$Z=Z_1+iZ_2$.  If $F$ is a smooth function on $G$ with
values in a complex vector space of finite dimension, define
\[
 (D_XF)(g)=\left.\frac{d}{dt}F(g\exp(tX))\right|_{t=0},
 \qquad E_{ab}=D_{(E^{ab})_1}+iD_{(E^{ab})_2},
\]
where $E^{ab}$ is the matrix unit in row $a$, column $b$.
The same definitions apply to smooth functions on an open subset
of $G$, by locality.  These are complex-linear differential
operators with smooth coefficients.  In particular, their
finite compositions preserve smoothness and do not enlarge
support.

For ordered injective lists $a=(a_0,\ldots,a_{m-1})$ and
$b=(b_0,\ldots,b_{m-1})$ of basis indices, define the column
determinant
\begin{equation}\label{ca-column}
 Q_{a,b}(u)=
 \sum_{\sigma\in S_m}\operatorname{sgn}(\sigma)
 \prod_{j=0}^{m-1}
 \bigl((u-j)\delta_{a_{\sigma(j)},b_j}\,\mathrm{Id}
                                  +E_{a_{\sigma(j)},b_j}\bigr).
\end{equation}
The product is in the displayed order; when it acts on a
function, the rightmost operator acts first.  No commutation
of differential operators is understood.
For $m=0$ the value is $\mathrm{Id}$.

\begin{lemma}\label{ca-column-leading}
The expression \eqref{ca-column} is a polynomial in $u$ of
degree at most $m$, with leading coefficient
\[
 [u^m]Q_{a,b}(u)=
 \det\bigl(\delta_{a_i,b_j}\bigr)_{0\leq i,j<m}\,\mathrm{Id}.
\]
This identity remains valid for arbitrary endomorphisms
$E_{ab}$ of any complex vector space; no Lie algebra
relations are needed.
\end{lemma}

\begin{proof}
The scalar $u$ commutes with every endomorphism.
On expanding each ordered product, a term of degree $m$
can only arise by choosing the scalar $u$ from every factor.
Its coefficient is
$\prod_j\delta_{a_{\sigma(j)},b_j}\,\mathrm{Id}$.
Summing over $\sigma$ gives the displayed determinant.
All other terms have degree at most $m-1$.
\end{proof}

For $0\leq m\leq r$ set
\[
 X_{a,c}(g)=(Hg)_{e_a,c},\qquad
 Y_{a,c}(g)=(g^{-1})_{c,e_a}\quad(0\leq a<m).
\]
There are two orders on the basis indices.  The first is the
displayed order used to define $H$.  The second begins with
$f_0,f_1,\ldots,f_{r-1}$ and is followed by
$e_0,\ldots,e_{r-1},w_0,\ldots,w_k$.
Write $c$ increasing with respect to the indicated order
in the following two sums:
\begin{align}
 (A_m(u)F)(g)
   &=\sum_{c\text{ increasing in the first order}}
          \det(X_{a,c_b}(g))_{a,b<m}\,
                 (Q_{(e_0,\ldots,e_{m-1}),c}(u)F)(g),
                    \label{ca-Aplus}\\
 (\widetilde A_m(u)F)(g)
   &=\sum_{c\text{ increasing in the second order}}
          \det(Y_{a,c_b}(g))_{a,b<m}\,
        (Q_{c,(f_0,\ldots,f_{m-1})}(u+m-1)F)(g).
                    \label{ca-Aminus}
\end{align}
Each determinant multiplies the result of the differential
operator; its entries are not differentiated in these formulas.
Put
\[
 d_m(g)=\det\bigl((Hg)_{e_a,e_b}\bigr)_{a,b<m},
 \qquad d_0(g)=1.
\]
These are the actual principal minors in the source at levels
$m\leq r$.  The additional vectors used for the final principal
minor do not occur at these levels.

\begin{proposition}\label{ca-A-leading}
For every smooth $F$ and every $g\in G$,
\[
 [u^m](A_m(u)F)(g)=d_m(g)F(g),\qquad
 [u^m](\widetilde A_m(u)F)(g)=\overline{d_m(g)}F(g).
\]
Both expressions have degree at most $m$.
\end{proposition}

\begin{proof}
In \eqref{ca-Aplus}, Lemma~\ref{ca-column-leading} gives a
nonzero leading coefficient only if the set of indices in
$c$ is $\{e_0,\ldots,e_{m-1}\}$.  The order condition then
forces $c=(e_0,\ldots,e_{m-1})$, with determinant of the
Kronecker matrix equal to one.  This proves the first formula.
The same argument in the second order forces
$c=(f_0,\ldots,f_{m-1})$.  The shift $u+m-1$ has no effect
on its leading coefficient.

Since $g^{-1}=Hg^*H$, one has
\[
 (g^{-1})_{f_b,e_a}
     =\overline{g_{f_a,e_b}}
     =\overline{(Hg)_{e_a,e_b}}.
\]
Taking determinants proves the second formula.  The empty
determinant case $m=0$ also gives both identities.
\end{proof}

Fix $0\leq j\leq r$ and work on the open set
$G_j=\{g\in G:d_j(g)\ne0\}$.
Let $c=(2r+k)/2$.  The two components of the operators
at a real parameter $y$ are
\begin{align}
 L^+_{j,e}(y)F&=d_j^{-1}A_{j-e}(y+c)F,
                    \label{ca-Lplus}\\
 L^-_{j,e}(y)F&=\overline{d_j}^{-1}
                       \widetilde A_{j-e}(y-c)F,
                    \label{ca-Lminus}
\end{align}
for $e=0,1$ and $e\leq j$.  Define
$L^\pm_{0,1}(y)F=0$.
These formulas include the special level-zero convention;
they never invoke an operator with a negative index.

\begin{corollary}\label{ca-normalized-leading}
As operators on smooth functions on $G_j$,
\[
 L^\pm_{j,0}(y)=y^j\mathrm{Id}
                       +\sum_{\ell<j}y^\ell B^\pm_\ell,\qquad
 L^\pm_{j,1}(y)=\sum_{\ell<j}y^\ell C^\pm_\ell,
\]
for differential operators with smooth coefficients independent
of $y$.  A translation $y\mapsto y+a$ by a fixed real number
does not change these leading coefficients.
At $j=0$, these formulas mean
$L^\pm_{0,0}=\mathrm{Id}$ and $L^\pm_{0,1}=0$.
\end{corollary}

\begin{proof}
Apply Proposition~\ref{ca-A-leading} with $m=j$ or $j-1$.
For $e=0$, multiplication by $d_j^{-1}$ or
$\overline{d_j}^{-1}$ cancels the leading coefficient
exactly.  For $e=1$ the degree is at most $j-1$.
The denominators are smooth and nonzero on $G_j$.
Translations preserve the leading coefficient of a polynomial.
\end{proof}

This proof also gives the needed regularity of every coefficient.
If a test function is smooth with compact support in $G_j$,
then each coefficient has support in the same compact set.
More generally, applying any fixed complex-linear restriction
or pullback map preserves the polynomial identity.  If a
Hermitian form is only defined on a subspace of functions,
all coefficient functions must belong to that subspace.
It suffices to check membership for $j+1$ distinct real
values of $y$: inversion of the scalar Vandermonde matrix
then expresses every coefficient as a finite linear combination
of those values.  This observation prevents any implicit
exchange of differentiation with an unproved integral.

\subsection{The leading coefficient in the mixed pairing}

In this subsection a sesquilinear form is linear in its first
argument and conjugate-linear in its second, which is the
convention of the source's coefficient functions.
Let $\mathcal B:V\times V'\to\mathbb C$ be such a form.
Suppose four vector polynomials in a real variable $y$ satisfy
\[
 X_0(y)=y^j x+\text{terms of degree less than }j,\quad
 Y_0(y)=y^j x'+\text{terms of degree less than }j,
\]
\[
 \deg X_1<j,\qquad \deg Y_1<j.
\]
For $j=0$ the last two polynomials are zero.
Let $U:V\to V$ be complex-linear and let $s\in\mathbb C$.
Then
\begin{equation}\label{ca-mixed-polynomial}
 F(y)=
 \mathcal B\bigl(U X_1(y)+sU X_0(y),Y_0(y)\bigr)
       -\mathcal B\bigl(U X_0(y),Y_1(y)\bigr)
\end{equation}
is a polynomial in the real variable $y$ of degree at most $2j$,
and
\begin{equation}\label{ca-mixed-leading}
 [y^{2j}]F(y)=s\,\mathcal B(Ux,x').
\end{equation}
Indeed, in the first term only
$\mathcal B(sy^j Ux,y^j x')$ contributes to that degree,
and the second term has degree at most $2j-1$.
Reality of $y$ is used when conjugating its powers.

For completeness the extra summand $sU X_0$ comes from an
explicit operation on the two-component indices, not from
an estimate.  Given a finite sequence of signs
$\varepsilon_\ell\in\{+,-\}$, regard its vectors as indexed by
$\{0,1\}^a$.  On the $\ell$th factor apply the matrix
\[
 \begin{pmatrix}1&0\\ s&1\end{pmatrix}
 \quad\text{if }\varepsilon_\ell=+,\qquad
 \begin{pmatrix}1&0\\ \overline s&1\end{pmatrix}
 \quad\text{if }\varepsilon_\ell=-.
\]
The matrices on distinct factors commute.  One may additionally
multiply every component by the same function
$\exp(sz-\overline s\,\overline z)$.
The latter has absolute value one, but its preservation of
a particular differential Hermitian form is an additional
identity, not a consequence of its pointwise absolute value.

On adding a new positive component, the transformed components
are precisely $UX_0$ and $UX_1+sUX_0$.
Contracting the new positive and negative components by
\[
 \mathcal B_1((x_0,x_1),(y_0,y_1))
           =\mathcal B(x_1,y_0)-\mathcal B(x_0,y_1)
\]
therefore gives \eqref{ca-mixed-polynomial}.
Repeated contraction gives a sesquilinear form on any number
of new components.
Applying Corollary~\ref{ca-normalized-leading} to a new pair
$L^+_{j,e}(y)$ and $L^-_{j,e}(y+1)$ proves
\eqref{ca-mixed-leading} for the actual normalized differential
operators.  The shift by one on the negative side changes no
leading coefficient.  No positivity is required for this
algebraic assertion.

\subsection{Exact extraction by interpolation}

For distinct real nodes $y_0,\ldots,y_d$, put
\[
 c_i=\prod_{\substack{0\leq h\leq d\\h\ne i}}
                    (y_i-y_h)^{-1},
 \qquad C=\sum_{i=0}^d|c_i|.
\]
Each $c_i$ is nonzero and $C>0$.
Lagrange interpolation of a vector polynomial $P$ of degree
at most $d$ gives
\begin{equation}\label{ca-Lagrange-leading}
 [y^d]P(y)=\sum_{i=0}^d c_i P(y_i).
\end{equation}
To prove this for any complex vector space, subtract the
usual interpolation expression from $P$.  Its coefficients
lie in a finite-dimensional subspace, so the scalar
root-counting argument applied to coordinates proves that
the difference vanishes.  Comparing coefficients of $y^d$
then gives \eqref{ca-Lagrange-leading}.
For $d=0$, it simply says that a constant polynomial equals
its value at its single node.

\begin{proposition}\label{ca-iterated-extraction}
Let $K,d$ be nonnegative integers.  Suppose that quantities
$W_a(y_1,\ldots,y_a)$, $0\leq a\leq K$, take values in a
complex vector space, with $W_0$ independent of the variables.
Assume that, for each fixed $y_1,\ldots,y_{a-1}$,
$W_a$ is a polynomial in its last real variable of degree
at most $d$, with leading coefficient $s_aW_{a-1}$.
For each $a$ choose $d+1$ distinct real nodes $y_{a,i}$ and
let $c_{a,i}$ be their coefficients above.  Then
\begin{equation}\label{ca-extraction}
 \left(\prod_{a=1}^K s_a\right)W_0
  =\sum_{i_1,\ldots,i_K}
       \left(\prod_{a=1}^K c_{a,i_a}\right)
           W_K(y_{1,i_1},\ldots,y_{K,i_K}).
\end{equation}
No assumption about simultaneous polynomial degree in all
variables is needed.
\end{proposition}

\begin{proof}
Extract the coefficient in the last variable using
\eqref{ca-Lagrange-leading}, then extract successively in
the preceding variables.  At each step the extracted
coefficient is the asserted scalar multiple of $W_{a-1}$,
independent of the variable just removed.
The finite sums commute with scalar multiplication.
The case $K=0$ has one empty tuple and is the identity.
\end{proof}

In the mixed pairing above one takes $d=2j$ and every
$s_a=s$.  Newly inserted operators may act on the previous
ones; they are not commuted past them.
Their leading term is the identity, so the same one-variable
argument applies to each successive insertion.
Thus \eqref{ca-extraction} proves the source's algebraic
degree and extraction identity on the common domain of its
forms and restrictions.

The natural probability weights are
\[
 \theta_{i_1,\ldots,i_K}
   =\frac{\prod_a|c_{a,i_a}|}{\prod_a C_a},
 \qquad C_a=\sum_i|c_{a,i}|.
\]
Their sum is one, because a product of finite sums expands
as a sum over all tuples.  Consequently, for vectors in a
normed space,
\begin{align}
 \left(\prod_a|s_a|\right)\|W_0\|
   &\leq\left(\prod_aC_a\right)
         \sum_{\boldsymbol i}\theta_{\boldsymbol i}
                    \|W_K(\boldsymbol y_{\boldsymbol i})\|,
                    \label{ca-extraction-norm}\\
 \left(\prod_a|s_a|\right)^2\|W_0\|^2
   &\leq\left(\prod_aC_a\right)^2
         \sum_{\boldsymbol i}\theta_{\boldsymbol i}
                    \|W_K(\boldsymbol y_{\boldsymbol i})\|^2.
                    \label{ca-extraction-square}
\end{align}
The first is the triangle inequality in
\eqref{ca-extraction}; the second follows from
Cauchy--Schwarz for the probability weights.

\begin{lemma}\label{ca-Chebyshev}
If $d\geq1$ and $\alpha<\beta$, the nodes
\[
 y_i=\frac{\alpha+\beta}{2}
       +\frac{\beta-\alpha}{2}\cos\frac{i\pi}{d},
       \qquad 0\leq i\leq d,
\]
satisfy the exact identity
\[
 \sum_i|c_i|=\frac12\left(\frac4{\beta-\alpha}\right)^d.
\]
This is the smallest possible coefficient sum among all
choices of $d+1$ distinct nodes in $[\alpha,\beta]$.
\end{lemma}

\begin{proof}
The polynomial $T_d$ determined by
$T_d(\cos\theta)=\cos(d\theta)$ has degree $d$, leading
coefficient $2^{d-1}$, and absolute value at most one on
$[-1,1]$.  These facts follow from
$T_0=1$, $T_1=x$, and $T_{d+1}=2xT_d-T_{d-1}$.
The scaled polynomial
\[
 R(y)=T_d\left(\frac{2y-\alpha-\beta}{\beta-\alpha}\right)
\]
has leading coefficient
$2^{2d-1}(\beta-\alpha)^{-d}$ and value $(-1)^i$ at $y_i$.
The nodes decrease with $i$, so $c_i$ also has sign $(-1)^i$.
Equation~\eqref{ca-Lagrange-leading} gives the equality.
For any other nodes in the interval, applying the same
identity to $R$ and using $|R|\leq1$ gives the asserted
lower bound.
\end{proof}

The actual construction needs its nodes separated from the
parameters already used.  The exact Chebyshev identity alone
does not supply those simultaneous separation conditions.
The source uses grids whose coefficient sums are bounded
by $(28N/I)^{2j}$.  Section~\ref{cg-section} supplies the
complete construction and improves $28$ to $24$.
For those actual grids,
\eqref{ca-extraction-norm}--\eqref{ca-extraction-square}
apply with that bound, with no loss from the number of grid
tuples.

\subsection{The explicit rational correction}

Let $q_0$ be a real-valued function, and let
$w=((\varepsilon_1,a_1),\ldots,(\varepsilon_L,a_L))$
be a finite sequence, with $\varepsilon_\ell\in\{+,-\}$
and $a_\ell\in\mathbb R$.  The modified function occurring
in the source is exactly
\begin{equation}\label{ca-qword}
 q_w(y)=q_0(y)
   \prod_{\varepsilon_\ell=+}
           \frac{y-a_\ell+1}{y-a_\ell}
   \prod_{\varepsilon_\ell=-}
           \frac{a_\ell-y}{a_\ell-y-1}.
\end{equation}
Only the displayed evaluation at $y$ is needed below.

\begin{lemma}\label{ca-q-distortion}
Suppose $D>1$, $q_0(y)\ne0$, and
\[
 |y-a_\ell|\geq D\quad(\varepsilon_\ell=+),\qquad
 |a_\ell-y-1|\geq D\quad(\varepsilon_\ell=-).
\]
Then $q_w(y)/q_0(y)>0$ and
\begin{equation}\label{ca-q-distortion-eq}
 \left|\log\frac{q_w(y)}{q_0(y)}\right|
                      \leq\frac{L}{D-1}.
\end{equation}
In particular, if $N\geq20$, $L\leq N$ and $D\geq N^2/8$,
then
\[
 e^{-20/49}\leq\frac{q_w(y)}{q_0(y)}\leq e^{20/49}.
\]
If also $q_0(y)>0$ and
$\sqrt{|P(y)P(y+1)|/q_0(y)}\leq A$, then
\begin{equation}\label{ca-node-majorant}
 q_w(y)>0,\qquad
 \sqrt{\frac{|P(y)P(y+1)|}{q_w(y)}}
                                      \leq e^{10/49}A.
\end{equation}
\end{lemma}

\begin{proof}
Every factor in \eqref{ca-qword} is $1+u$ with
$|u|\leq1/D<1$.  It is positive, and
\[
 |\log(1+u)|
  =\left|\int_0^u\frac{dt}{1+t}\right|
  \leq\frac{|u|}{1-|u|}\leq\frac1{D-1}.
\]
Adding these inequalities proves
\eqref{ca-q-distortion-eq}.
Moreover
\[
 \frac{L}{D-1}\leq\frac{8N}{N^2-8}
                          \leq\frac{20}{49}\quad(N\geq20).
\]
The last function decreases for $N>\sqrt8$; at $N=20$
its value is $20/49$.
Taking exponentials gives the first uniform assertion,
and taking the reciprocal square root gives
\eqref{ca-node-majorant}.
\end{proof}

In the actual construction $N=20(r+1)$.
The separation conditions proved for its new and previous
parameters give precisely the denominator bounds in this
lemma, with $D=N^2/8$ in the exceptional short interval
and a larger value in the other intervals.
The polynomial $q_0$ is $Q$ from Section~\ref{as-section}.
On an interval belonging to a differing binary position it
is positive: exactly one of the paired factors in $Q$ has
opposite signs, and its defining initial minus sign reverses
that sign.  Thus \eqref{ca-node-majorant} applies there.
The source uses the weaker factor $eA$.
This improvement concerns the rational modification alone;
it does not assert that the separated grids or positive
Hermitian forms exist for arbitrary parameters.

\subsection{The exact positive matrix estimate}

We next give the finite-dimensional step that turns the
interpolation identity into the normalized local coefficient
estimate.  All its hypotheses are stated for arbitrary
linear combinations, as required by positivity of a block
matrix.

Fix $K,d\geq0$ and the nodes of
Proposition~\ref{ca-iterated-extraction}.
For a tuple $\boldsymbol i=(i_1,\ldots,i_K)$ put
\[
 c_{\boldsymbol i}=\prod_a c_{a,i_a},\quad
 C_*=\prod_a\sum_i|c_{a,i}|,\quad
 \theta_{\boldsymbol i}=|c_{\boldsymbol i}|/C_*.
\]
Let $I_+$ and $I_-$ be finite sets.  For each tuple suppose
$R^+_{\boldsymbol i}$ and $R^-_{\boldsymbol i}$ are
positive semidefinite Hermitian forms on
$\mathbb C^{I_+}$ and $\mathbb C^{I_-}$, and that
$B_{\boldsymbol i}$ is a sesquilinear form between these
two spaces.  Let $P$ and $q$ be real-valued at the nodes,
with $P(y_{a,i})P(y_{a,i}+1)\ne0$ and $q(y_{a,i})>0$.
Assume
\begin{equation}\label{ca-pairing-assumption}
 |B_{\boldsymbol i}(x,y)|^2
       \prod_{a=1}^Kq(y_{a,i_a})
       \leq R^+_{\boldsymbol i}(x,x)
                 R^-_{\boldsymbol i}(y,y)
 \quad\text{for every }x,y.
\end{equation}
Fix positive $\Gamma,S$ and nonnegative $T,A$ such that
\[
 \frac{|P(y_{a,i})P(y_{a,i}+1)|}{q(y_{a,i})}
                                                  \leq A^2.
\]
Define normalized forms
\[
 G^+_{\boldsymbol i}
       =\frac{R^+_{\boldsymbol i}}
                 {\Gamma S\prod_a|P(y_{a,i_a})|},\qquad
 G^-_{\boldsymbol i}
       =\frac{R^-_{\boldsymbol i}}
                 {\Gamma S\prod_a|P(y_{a,i_a}+1)|},
\]
and the mixed form
\[
 F=\frac{T^K}{\Gamma S}
            \sum_{\boldsymbol i}c_{\boldsymbol i}B_{\boldsymbol i}.
\]
Write $\overline G^\pm=\sum_{\boldsymbol i}
\theta_{\boldsymbol i}G^\pm_{\boldsymbol i}$.

\begin{proposition}\label{ca-positive-block}
For all $x\in\mathbb C^{I_+}$ and $y\in\mathbb C^{I_-}$,
\begin{equation}\label{ca-positive-bound}
 |F(x,y)|\leq C_*T^KA^K
       \sqrt{\overline G^+(x,x)\,\overline G^-(y,y)}.
\end{equation}
Consequently, if $X\geq0$ and
$C_*T^KA^K\leq X^K$, the block Hermitian form with
diagonal forms $X^K\overline G^+$, $X^K\overline G^-$
and mixed form $F$ is positive semidefinite.
\end{proposition}

\begin{proof}
Divide \eqref{ca-pairing-assumption} by $(\Gamma S)^2$.
The normalizations give exactly
\[
 \left|\frac{B_{\boldsymbol i}(x,y)}{\Gamma S}\right|
 \leq
 \left(\prod_a
       \frac{|P(y_{a,i_a})P(y_{a,i_a}+1)|}{q(y_{a,i_a})}
       \right)^{1/2}
       \sqrt{G^+_{\boldsymbol i}(x,x)
                         G^-_{\boldsymbol i}(y,y)}.
\]
In particular the factor on the right is at most $A^K$.
The triangle inequality and
$|c_{\boldsymbol i}|=C_*\theta_{\boldsymbol i}$ yield
\[
 |F(x,y)|\leq C_*T^KA^K
     \sum_{\boldsymbol i}\theta_{\boldsymbol i}
          \sqrt{G^+_{\boldsymbol i}(x,x)
                            G^-_{\boldsymbol i}(y,y)}.
\]
Cauchy--Schwarz for the probability weights proves
\eqref{ca-positive-bound}.
If $a=\overline G^+(x,x)$ and
$b=\overline G^-(y,y)$, the quadratic value of the block
form is at least
$X^K(a+b)-2|F(x,y)|\geq X^K(\sqrt a-\sqrt b)^2$.
This is nonnegative, including when one of the diagonal
values is zero.  For $K=0$ all empty products are one and
the same proof applies.
\end{proof}

In the source the original size $S$ is the product
$\prod_{\ell\in w}|P(a_\ell)|$.
Adding positive parameters $y_1,\ldots,y_K$ multiplies it by
$\prod_a|P(y_a)|$; adding the corresponding negative
parameters $y_a+1$ multiplies it by
$\prod_a|P(y_a+1)|$.  Thus the displayed normalization is
exactly the one in its Gram sums.  The factor $\Gamma$
cancels in the estimate.  It has not been estimated or
replaced by one.

For the Fourier coefficient at level $j$, the interpolation
degree is $d=2j$.  If each actual grid has
$C_a\leq(28N/I)^{2j}$, and $A_{\rm pol}$ is the unscaled
boundary majorant
$\sup\sqrt{|P(y)P(y+1)|/q_0(y)}$ on the relevant interval,
one may take
\begin{equation}\label{ca-X-exact}
 X=\frac{g\lambda}{2\pi}
         \left(\frac{28N}{I}\right)^{2j}
                         e^{10/49}A_{\rm pol}.
\end{equation}
Here $g>0$ is the level parameter and $\lambda>0$ is the
actual scalar change of the complex integration variable.
This conclusion uses \eqref{ca-pairing-assumption} for the
particular Hermitian forms.  Its proof has not supplied
that representation-theoretic hypothesis.

\subsection{Fourier normalization and the remaining analytic input}

The source's additive character at a complex place is
\[
 \psi_g(l,c)=\exp\left(\frac{4\pi i}{g}\Re(lc)\right).
\]
If the radial variable is $z=\lambda c$, put
\[
 s=\frac{2\pi i\,l}{g\lambda},\qquad
 T=\frac{g\lambda}{2\pi}.
\]
Then the exact phase and scalar identities are
\begin{equation}\label{ca-phase}
 e^{sz-\overline s\,\overline z}=\psi_g(l,c),
 \qquad T^Ks^K=i^Kl^K.
\end{equation}
These follow by substituting $z=\lambda c$ and using that
$g,\lambda$ are real.  In the usual Fourier transform
$\widehat f(\xi)=\int_{\mathbb C}f(c)
e^{-2\pi i\Re(\overline\xi c)}\,dc$, the same character
therefore has frequency $\xi=-2\overline l/g$.

If the actual radial Fourier identity represents its
normalized coefficient as
\[
 \widehat K(l)
     =(\Gamma S)^{-1}\,
               \mathcal B(U_sx,x'),
\]
then the proved extraction identity identifies the mixed
form in Proposition~\ref{ca-positive-block} with
$i^Kl^K\widehat K(l)$.
Here $U_s$ consists of the explicit phase and the component
matrices described above.  There is a second necessary
identity: $U_s$ must preserve the original nonnegative
Hermitian form.  If it does, Cauchy--Schwarz supplies the
coefficient bound with no power of $l$ as well.
The diagonal forms in this latter bound need not be the
averaged diagonal forms in Proposition~\ref{ca-positive-block}.
If both diagonal products are bounded by a common
nonnegative majorant $M^2$, taking the stronger estimate
gives
\[
 |\widehat K(l)|\leq
       M\min\bigl(1,(X/|l|)^K\bigr)\qquad(l\ne0).
\]
At $l=0$ one uses the bound $M$ from the original form.
Obtaining this common majorant is a separate requirement
when assembling the global estimate.

Thus the differential-operator leading terms, real polynomial
extraction, rational correction, finite averaging, and exact
normalizations have now been proved.  The remaining actual
inputs are the radial Fourier identity, nonnegativity and
the mixed inequality \eqref{ca-pairing-assumption} for the
specified differential Hermitian forms, preservation of those
forms by $U_s$, and the analytic function-domain conditions.
The finite grid and separation conditions are discharged in
Section~\ref{cg-section}.  The mass,
source-dimension, and formal-degree hypotheses in
Section~\ref{as-section} also remain separate.

\begin{remark}\label{ca-counterexamples}
Two possible shortcuts are false.  First, entrywise bounds
do not imply the block positivity used here.  Let the two
diagonal matrices be $I_2$ and the mixed matrix have every
entry equal to one.  Each mixed entry has absolute value
at most the geometric mean of its two diagonal entries.
For $x=(1,1)$ the block quadratic form at $(x,-x)$ is
$4-8=-4$.  This is why \eqref{ca-pairing-assumption} is
required for all linear combinations.

Second, the real parameter in
\eqref{ca-mixed-polynomial} cannot be replaced by an
arbitrary complex parameter while retaining the same
polynomial identity.  With the ordinary form on
$\mathbb C$, $X_0(y)=Y_0(y)=y$ gives
$\mathcal B(X_0(y),Y_0(y))=|y|^2$, not $y^2$.
At $y=i$ these values are $1$ and $-1$.
The actual source uses real interpolation nodes and
tests the mixed form on all linear combinations, so
neither example contradicts its stated hypotheses.
\end{remark}


